% =====================================================================
%  IEEE-STYLE CONFERENCE REPORT  (self-contained)
%  MLE-AI-201 Machine Learning -- Final Project
%
%  Build:   pdflatex report.tex ; pdflatex report.tex      (as CI does)
%  XeTeX/LuaTeX engines (tectonic, xelatex, lualatex) would drop the
%  Azerbaijani letters in the generated tables, so the build stops there.
%
%  Tables and figures marked "generated" are written by
%      python -m src.run_all
%  into report/tables/ and report/figures/ -- never edited by hand.
% =====================================================================
\documentclass[10pt,twocolumn,a4paper]{article}

% ---- Geometry: IEEE-like margins / two narrow columns ----------------
\usepackage[a4paper,top=0.75in,bottom=1.0in,left=0.68in,right=0.68in,columnsep=0.25in]{geometry}

% ---- Times-like font (standard, ships with every TeX distro) ---------
\usepackage{mathptmx}          % Times text + math
\usepackage[T1]{fontenc}
\usepackage[utf8]{inputenc}    % allows accented / Azerbaijani chars

% ---- Common, dependency-free packages --------------------------------
\usepackage{graphicx}
\usepackage{amsmath,amssymb}
\usepackage{booktabs}
\usepackage{array}
\usepackage{titlesec}
\usepackage{caption}
\usepackage{cite}
\usepackage[hidelinks]{hyperref}
\usepackage{url}

% ---- IEEE-style section headings -------------------------------------
\renewcommand{\thesection}{\Roman{section}}
\renewcommand{\thesubsection}{\Alph{subsection}}
\titleformat{\section}[block]{\centering\normalsize\scshape}{\thesection.}{0.5em}{}
\titleformat{\subsection}[block]{\normalsize\itshape}{\thesubsection.}{0.5em}{}
\titlespacing*{\section}{0pt}{1.4ex plus 1ex minus .2ex}{1ex}
\titlespacing*{\subsection}{0pt}{1.0ex plus 1ex minus .2ex}{0.6ex}

% ---- Abstract / Index-Terms in the IEEE boldface-dash style ----------
\newcommand{\abstractkw}[1]{\noindent\textbf{\textit{#1}}}

\captionsetup{font=small,labelsep=period}
\setlength{\parindent}{1em}

% ---- Azerbaijani names in the generated tables (districts, categories):
% with pdflatex the T1 font has every letter but the schwa, so draw it as a
% turned e. XeTeX/LuaTeX would silently drop several letters with this font.
\ifdefined\DeclareUnicodeCharacter
  \DeclareUnicodeCharacter{0259}{\rotatebox[origin=c]{180}{e}}
  \DeclareUnicodeCharacter{018F}{\scalebox{1.3}{\rotatebox[origin=c]{180}{e}}}
\else
  \errmessage{Build report.tex with pdflatex, as CI does: this engine would
    drop the Azerbaijani letters in the generated tables}
\fi

% ---- Tables and figures written by `python -m src.run_all` -----------
% (report/tables/, report/figures/). Until they are generated, a box
% names the missing file instead of breaking the build.
\newcommand{\missinggenerated}[1]{\fbox{\parbox{0.95\linewidth}{\small
  Missing \nolinkurl{report/#1}: run \texttt{python -m src.run\_all}.}}}
\newcommand{\generatedtable}[1]{\IfFileExists{#1}{\input{#1}}{%
  \par\noindent\missinggenerated{#1}\par}}
\newcommand{\generatedfigure}[2]{\IfFileExists{#1}{%
  \includegraphics[width=#2]{#1}}{\missinggenerated{#1}}}

% =====================================================================
\begin{document}

% ------------------------- TITLE BLOCK --------------------------------
\twocolumn[
  \begin{@twocolumnfalse}
  \begin{center}
    {\LARGE\bfseries Pricing bina.az Listings from Scratch:\\A NumPy Decision Tree and a Pegasos SVM
      on 63{,}463 Azerbaijani Property Listings\par}
    \vspace{1.2em}
    % Authors: Surname, Given name --- alphabetical by surname.
    {\large
      Eyyubov, Samur\quad Huseynli, Ali\quad Ibrahimli, Polad\quad Usubov, Tural\par}
    \vspace{0.5em}
    {\itshape MLE-AI-201 Machine Learning, AI Academy --- Cohort II\par}
    \vspace{1.4em}
  \end{center}
  \begin{quote}
    \small
    \abstractkw{Abstract---}% IEEE uses an em-dash after "Abstract"
    We predict the asking price of property listings scraped from
    \emph{bina.az}, Azerbaijan's largest real-estate platform, with two models
    written from first principles in NumPy: a CART decision tree
    (Gini/entropy/MSE) and a linear soft-margin SVM trained by the Pegasos
    sub-gradient method. A documented cleaning pipeline turns a raw dump of
    100{,}775 scrape rows into 63{,}463 unique listings and 115 encoded
    features, drops every column derived from the target, and produces a
    seeded 70/15/15 split. One dataset carries two tasks: regression on
    $\log(\text{price})$ and premium-vs-standard classification at the
    training-median price. Hyperparameters are chosen on the validation split
    only, and every model is compared against scikit-learn with matched
    settings, five-fold cross-validation, 95\% bootstrap intervals and
    wall-clock timings. Our tree reaches a test RMSE of $0.332$ on
    $\log(\text{price})$ ($R^2=0.829$, median error $10.3\,\%$ in AZN) ---
    the same to three decimals as \texttt{DecisionTreeRegressor}, with
    $99.2\,\%$ of its test predictions identical --- and classifies the price
    tier with F1 $0.921$ against scikit-learn's $0.919$, while our Pegasos SVM
    scores F1 $0.867$, above \texttt{SGDClassifier} ($0.849$) and below the
    converged \texttt{LinearSVC} ($0.893$). Writing the models ourselves costs
    a constant factor of 13--40 in fit time, not accuracy; a bonus forest of
    our own trees reaches RMSE $0.291$ and F1 $0.944$. Every number and figure
    here comes from one seeded command, \texttt{python -m src.run\_all}:
    \href{https://github.com/eyyubovsamur240-afk/ML-final-project-main}{\texttt{github.com/eyyubovsamur240-afk/ML-final-project-main}}
    (tag \texttt{v1.0-final}).

    \vspace{0.6em}
    \abstractkw{Index Terms---}decision tree, CART, support-vector machine,
    Pegasos, sub-gradient descent, regression, binary classification,
    real-estate price prediction, reproducibility.
  \end{quote}
  \vspace{1.0em}
  \end{@twocolumnfalse}
]

% =====================================================================
\section{Introduction}
Asking prices on \emph{bina.az} are set by owners and agencies, not by a
market maker, and the platform publishes them next to a handful of structured
attributes: floor area, room count, floor, repair state, coordinates and a
district name. Learning the mapping from those attributes to the price is a
useful exercise in two ways at once. It is a real pricing problem --- the same
estimate a buyer wants before making an offer --- and the data is a genuine
scraped dump, with mixed types, Azerbaijani column names, repeated scrapes,
re-posted listings, structurally missing fields and data-entry errors that
run from an 11~AZN commercial unit to a 600~M~AZN flat.

This project builds the two models from scratch rather than calling a library,
and treats scikit-learn strictly as the thing to be compared against. Our
contributions are:

\begin{enumerate}
  \item \textbf{A documented cleaning pipeline} (\texttt{src/data\_prep.py})
    that turns 100{,}775 raw rows into 63{,}463 unique listings and 115
    encoded features, removes every column that is a function of the target,
    and produces a seeded, leakage-free 70/15/15 split whose price-tier
    threshold is fixed on the training rows only.
  \item \textbf{A CART decision tree in NumPy} (\texttt{src/decision\_tree.py})
    supporting Gini, entropy and MSE, with a vectorised prefix-sum split
    search, scikit-learn's stopping-rule semantics, \texttt{predict\_proba}
    from leaf frequencies and impurity-decrease importances.
  \item \textbf{A Pegasos soft-margin SVM} (\texttt{src/svm.py}) that
    minimises the regularised hinge-loss primal by stochastic sub-gradient
    descent and records its objective while it trains.
  \item \textbf{A matched benchmark} (\texttt{src/baselines.py},
    \texttt{src/tree\_experiments.py}) in which scikit-learn's trees get our
    tuned hyperparameters and our rows, and \texttt{SGDClassifier} and
    \texttt{LinearSVC} are given the same $\lambda$-regularised hinge
    objective, down to the number of updates.
  \item \textbf{An evaluation and diagnosis layer} (\texttt{src/evaluate.py},
    \texttt{src/error\_analysis.py}) in which every metric, the stratified
    $k$-fold splitter and the bootstrap are written from scratch, and the test
    errors are broken down by price band, property category, district coverage
    and distance from the tier threshold.
  \item \textbf{Bonus Problem~9:} a random forest grown from our own trees,
    with the test score tracked as learners are added.
  \item \textbf{One-command reproduction} (\texttt{src/run\_all.py}): a single
    seeded entry point writes every number, table and figure used here, and a
    test asserts that two runs agree except for the timings.
\end{enumerate}

\section{Related Work and Background}
\textbf{Decision trees.} We implement CART~\cite{cart}: a greedy, binary,
axis-aligned recursive partition. A node holding rows $S$ with impurity $I(S)$
is split at feature $j$ and threshold $\tau$ into $S_L=\{x: x_j\le\tau\}$ and
$S_R$, and the split is scored by the impurity decrease
\begin{equation}
  \mathrm{IG} = I(S) - \frac{|S_L|}{|S|}I(S_L) - \frac{|S_R|}{|S|}I(S_R),
\end{equation}
with $I$ the Gini index $1-\sum_k p_k^2$, the entropy $-\sum_k p_k\log_2 p_k$
or, for regression, the variance $\frac{1}{|S|}\sum_i (y_i-\bar y)^2$. The
greedy search is exact for a single node and heuristic for the tree as a
whole; depth and leaf-size limits are the standard way of trading bias for
variance. scikit-learn's \texttt{DecisionTreeClassifier} and
\texttt{DecisionTreeRegressor}~\cite{sklearn} implement the same algorithm,
which is what makes a node-by-node comparison meaningful.

\textbf{Support-vector machines.} The soft-margin SVM~\cite{cortes} separates
two classes with the largest-margin hyperplane while paying a hinge penalty
for violations. Written as an unconstrained primal, it minimises
$\frac{\lambda}{2}\lVert\mathbf{w}\rVert^2 + \frac1n\sum_i\max(0,1-y_i f(x_i))$,
a convex but non-differentiable objective. Rather than solving the dual
quadratic program, we use Pegasos~\cite{pegasos}, which takes a stochastic
sub-gradient step on one sampled example per iteration with the decaying rate
$\eta_t = 1/(\lambda t)$; it reaches a fixed accuracy in a number of
iterations that does not grow with the training-set size, which is exactly why
it is a practical choice for a from-scratch implementation. Our scikit-learn
twins, \texttt{SGDClassifier} with the hinge loss and \texttt{LinearSVC},
minimise the same objective with a different optimiser and a different
intercept convention; Section~\ref{sec:setup} says how we matched them.

\section{Data and Preprocessing}
\label{sec:data}
\textbf{The raw dump.} The \emph{Bina.az Sale Project} dataset (Kaggle,
\texttt{sehriyarmemmedli/binaaz-sale-project})~\cite{dataset} contains 100{,}775
rows and 51 columns scraped between October 2024 and February 2025. Every
price is quoted in AZN (both currency columns are constant), so no currency
conversion is needed, but almost everything else needs work: numeric fields
arrive as text with units (\texttt{Sahə} = ``145~m\textsuperscript{2}'',
``1.3 sot''), the floor is a string (``7 / 9''), the column names and values
are Azerbaijani, and the same listing appears many times because the site was
scraped repeatedly.

\textbf{Cleaning.} Table~\ref{tab:cleaning} is the ledger of what we removed
and why; the code is \texttt{data\_prep.clean\_data}, and the notebook
\texttt{eda.ipynb} shows the evidence behind each rule.

\begin{table}[t]
  \centering
  \caption{Cleaning ledger (\texttt{src/data\_prep.py}). Rules are applied in
    this order; the counts are what each step removed.}
  \label{tab:cleaning}
  \footnotesize
  \begin{tabular}{@{}lrr@{}}
    \toprule
    Step & Removed & Rows left \\
    \midrule
    Raw dump                                      & 0      & 100{,}775 \\
    Repeated scrapes of the same URL              & 36{,}321 & 64{,}454 \\
    Outliers (fixed domain rules)                 & 81     & 64{,}373 \\
    Same listing ID, different URL format         & 669    & 63{,}704 \\
    Exact re-posts (same property and description)& 53     & 63{,}651 \\
    Houses with an unknown room count             & 187    & 63{,}464 \\
    Coordinates outside Azerbaijan                & 1      & 63{,}463 \\
    \bottomrule
  \end{tabular}
\end{table}

\emph{Repeated scrapes.} 20{,}506 listings appear more than once (up to 17
copies) and 3{,}356 of them changed price between scrapes, so we keep the most
recent scrape of each listing URL rather than the first. Two further
de-duplication passes catch what the URL key misses: 669 listings whose URL
differs only by a \texttt{.html} suffix (keyed on the numeric listing ID), and
53 exact re-posts of the same property under a new ID, matched on category,
price, areas, rooms, floor, \emph{exact} coordinates and an identical
description. The key is deliberately strict: identical flats in one new
building are genuinely different listings and are kept.

\emph{Outliers.} We use fixed, category-aware domain rules, never a quantile
learned from the data: buildings under 5{,}000~AZN, residential listings under
15{,}000~AZN or below 15~m\textsuperscript{2}, more than
300~m\textsuperscript{2} per room, a non-commercial price per
m\textsuperscript{2} outside 100--15{,}000~AZN (commercial: 100--60{,}000),
land outside 0--10{,}000 sot, and land under 50~AZN per sot. Together they
remove 81 rows (0.13\,\%). Price per m\textsuperscript{2} and price per sot
are computed \emph{only} to flag these rows and never enter a feature matrix.

\emph{Types and missing values.} Areas are parsed into a floor area in
m\textsuperscript{2} and a land area in sot (land plots store their size in
sot in the area column and have no building, so their floor area is 0);
``7 / 9'' becomes floor and total floors; \texttt{Təmir}
(var/yoxdur) becomes yes/no/unknown. Values that are missing \emph{by design}
are filled with the meaning they carry --- no separate land for flats, no
floor for houses, land and commercial units, no room count for land, objects
and garages --- and the room count is first recovered from the free-text
attribute string (``4 otaqlı, 145~m\textsuperscript{2}''). Only genuinely
unknown values are dropped: 187 houses with no recoverable room count. After
cleaning the modelling table has no missing values at all, which the pipeline
asserts.

\textbf{Leakage columns removed.} \texttt{unit\_price} (price per
m\textsuperscript{2}) and \texttt{total\_price} (the price restated) are
arithmetic functions of the target and are dropped; a test and a runtime
assertion keep them out of the clean table. We also drop identifiers
(\texttt{owner\_name}, \texttt{shop\_name}, \texttt{address}, URLs and listing
IDs are kept only for de-duplication and bookkeeping), \texttt{views} (known
only after a listing has been live), the duplicated Latin copies of three
Azerbaijani columns, \texttt{Binanın növü} (99.8\,\% missing), and the
dates, image URLs and free text.

\textbf{Features.} We keep eight numeric features (floor area, land area,
rooms, floor, total floors, latitude, longitude and the haversine distance to
Fountains Square, the Baku city centre), four binary flags (bill of sale,
mortgage, VIP, featured) and three categoricals (property category, repair
state, district), one-hot encoded. Districts with fewer than 30 \emph{training}
listings are grouped into ``other'' (92 of 142 districts survive that filter),
which keeps the encoding from exploding into near-empty columns: 115 feature
columns in total. The encoder is fitted on the training split alone, and a
category never seen in training becomes an all-zero row of its one-hot block.

\textbf{The two tasks and the split.} Task~A is regression on
$\log(\text{price})$: prices are extremely heavy-tailed (skewness 19.8, median
210{,}000~AZN, mean 327{,}199~AZN, 583 listings above the 99th percentile of
2.5~M~AZN), while $\log(\text{price})$ is nearly symmetric (skewness 0.20) ---
Figure~\ref{fig:price}. Task~B is binary classification into ``premium'' and
``standard'' at the \emph{training} median price, 209{,}000~AZN. The split
(seed 42) first cuts 70\,\% training and 30\,\% holdout stratified by property
category, then fixes the threshold on the training prices only, and finally
splits the holdout into validation and test stratified by the resulting tier
label. The class balance is therefore 50.0\,\% premium in training and
50.4\,\% in validation and test (Figure~\ref{fig:balance}). Table~\ref{tab:data}
summarises the modelling dataset. No
preprocessing statistic --- the threshold, the district whitelist, the one-hot
layout, the standardisation mean and std --- is ever computed on validation or
test rows.

\begin{table}[t]
  \centering
  \caption{Dataset summary after cleaning.}
  \label{tab:data}
  \begin{tabular}{@{}lr@{}}
    \toprule
    Property & Value \\
    \midrule
    Rows (raw dump)             & 100{,}775 \\
    Rows (after cleaning)       & 63{,}463 \\
    Features used (one-hot encoded) & 115 \\
    Train / Val / Test split    & 44{,}423 / 9{,}520 / 9{,}520 \\
    Price-tier threshold (train median) & 209{,}000 AZN \\
    Price-tier balance, premium (\%) & 50.0 / 50.4 / 50.4 \\
    \bottomrule
  \end{tabular}
\end{table}

\textbf{What the data looks like.} Floor area is the strongest single
predictor (Spearman $\rho=0.74$ against price, Figure~\ref{fig:size-price}a),
the room count follows it ($\rho=0.52$, Figure~\ref{fig:size-price}b), and
location matters through both distance to the centre ($\rho=-0.40$) and the
district itself: the median price of Ağ şəhər
(539{,}594~AZN) is more than eight times that of Binə q.
(65{,}000~AZN). Figure~\ref{fig:location} maps both effects.

% Figures 1-4 come from src/eda.py (run by run_all); their image titles carry
% the same numbers, so they stay the first four figures in the source.
\begin{figure*}[t]
  \centering
  \generatedfigure{figures/fig1_price_distribution.png}{0.8\textwidth}
  \caption{Sale price of the cleaned listings: (a) in thousand AZN, with the
    axis cut at the 99th percentile, and (b) $\log(\text{price})$, the scale
    the Task~A models are trained on.}
  \label{fig:price}
\end{figure*}

\begin{figure*}[t]
  \centering
  \generatedfigure{figures/fig2_area_rooms_vs_price.png}{0.8\textwidth}
  \caption{(a) Floor area against price (log--log) and (b) price by number
    of rooms, with Spearman correlations; land plots are left out because
    they have no floor area.}
  \label{fig:size-price}
\end{figure*}

\begin{figure*}[t]
  \centering
  \generatedfigure{figures/fig3_location_effects.png}{0.8\textwidth}
  \caption{Location and price: (a) every listing coloured by
    $\log(\text{price})$, (b) the median $\log(\text{price})$ per map cell in
    the Baku area.}
  \label{fig:location}
\end{figure*}

\begin{figure}[t]
  \centering
  \generatedfigure{figures/fig4_class_balance.png}{\linewidth}
  \caption{Share of standard and premium listings in the training,
    validation and test splits; the threshold is the training median price.}
  \label{fig:balance}
\end{figure}

\section{Method}
\subsection{Decision tree (from scratch)}
\texttt{src/decision\_tree.py} is a recursive CART implementation in NumPy
that serves both tasks: \texttt{criterion} is \texttt{"gini"} or
\texttt{"entropy"} for classification and \texttt{"mse"} for regression, and
the class exposes scikit-learn's API (\texttt{fit}, \texttt{predict},
\texttt{predict\_proba}, \texttt{apply}, \texttt{get\_depth},
\texttt{get\_n\_leaves}, \texttt{feature\_importances\_}), which is what lets
the benchmark swap one model for the other.

\textbf{Vectorised split search.} The naive search --- for every feature, for
every candidate threshold, recompute both children --- is $O(n^2d)$ per node
and far too slow at $n=44{,}423$ rows and $d=115$ columns. Instead, at each
node we sort \emph{all} feature columns at once
(\texttt{np.argsort(\dots, axis=0)}) and reorder the labels accordingly. A
cumulative sum down each sorted column then gives the left child's statistics
for \emph{every} split position simultaneously, and the right child's as the
total minus the prefix: class counts per class for Gini and entropy, and
$\sum y$ and $\sum y^2$ (after centring, to keep the subtraction numerically
stable) for the variance criterion. One $(n-1)\times d$ array therefore holds
the weighted child impurity of every candidate split, and the best split is a
single \texttt{argmin}. The cost per node is $O(nd\log n)$ for the sort and
$O(nd)$ (or $O(ndK)$ for $K$ classes) for the scores, with no Python loop over
thresholds; a Python loop over features alone made fully grown trees about six
times slower in our own measurements.

\textbf{Thresholds and determinism.} A split position is valid only where
consecutive sorted values differ by more than $10^{-7}$ (scikit-learn's
constant) and both children keep at least \texttt{min\_samples\_leaf} rows.
The threshold is the midpoint of the two neighbouring values, falling back to
the lower value when the midpoint rounds up to the upper one. Ties between
equally good splits are broken deterministically --- lowest feature index,
then smallest threshold --- whereas scikit-learn visits features in a random
order; this is the main reason two correct trees can differ on real data, and
Section~\ref{sec:results} quantifies it.

\textbf{Growth and stopping.} All four stopping rules required by the brief
are implemented with scikit-learn's semantics, so that matched hyperparameters
mean the same thing in the benchmark: \texttt{max\_depth} (never split at that
depth), \texttt{min\_samples\_split}, \texttt{min\_samples\_leaf} (enforced
inside the split search, not after it) and \texttt{min\_impurity\_decrease},
compared against the decrease weighted by the node's share of the training
rows, $\frac{N_t}{N}\,\mathrm{IG}$. A node is also a leaf once it is pure.
Leaves store class frequencies (classification) or the mean (regression), so
\texttt{predict} returns the majority class or the mean and
\texttt{predict\_proba} returns the leaf frequencies. Prediction pushes whole
index arrays down the tree instead of walking row by row, and
\texttt{feature\_importances\_} accumulates $N_t \cdot \mathrm{IG}$ per
feature, normalised to sum to one.

\textbf{Tests.} \texttt{tests/test\_decision\_tree.py} checks the tree against
hand-computed values (a root split whose gain is exactly $1/3$) and against
scikit-learn on tiny, \emph{fully determined} examples --- examples where
scikit-learn builds the same tree for 20 different random seeds, so no tie can
hide a bug --- for Gini, entropy and MSE, and on random data it asserts that
whenever scikit-learn itself is tie-free our tree puts exactly the same
training rows in the same leaves. Four further tests check that each stopping
rule actually changes the tree.

\subsection{Support-vector machine (from scratch)}
\texttt{src/svm.py} implements the linear soft-margin SVM as plain
sub-gradient descent on the regularised hinge-loss primal (no QP, no SMO):
\begin{equation}
  J(\mathbf{w},b)=\tfrac{\lambda}{2}\lVert\mathbf{w}\rVert^{2}
  +\tfrac{1}{n}\sum_{i=1}^{n}\max\!\big(0,\,1-y_i(\mathbf{w}^{\top}\mathbf{x}_i+b)\big),
\end{equation}
with labels mapped to $y_i\in\{-1,+1\}$ (the original labels are stored and
restored by \texttt{predict}). We fold the bias into the weight vector by
augmenting every row with a constant 1, so one update handles $\mathbf{w}$ and
$b$ together. At step $t$ Pegasos~\cite{pegasos} draws one training row $i$
uniformly with replacement, sets $\eta_t = 1/(\lambda t)$, and applies
\begin{equation}
  \mathbf{w} \leftarrow (1-\eta_t\lambda)\,\mathbf{w}
  + \eta_t\, y_i \mathbf{x}_i \cdot \mathbf{1}[\,y_i\mathbf{w}^{\top}\mathbf{x}_i < 1\,].
\end{equation}
Every step shrinks the weights; a step on a margin violator also pulls the
hyperplane towards that example. We do not use the optional projection onto
the ball of radius $1/\sqrt{\lambda}$ from the original paper. Because the
bias travels inside $\mathbf{w}$ it is regularised too, which is a small,
deliberate difference from \texttt{SGDClassifier} and \texttt{LinearSVC},
whose intercepts are not.

\textbf{Scaling and convergence.} The hinge loss is scale-sensitive, so every
SVM --- ours and both baselines --- sees features standardised with the
statistics of its own fit rows (\texttt{baselines.StandardScaled}, and
\texttt{StandardScaler} inside the scikit-learn pipelines), which keeps
cross-validation folds from leaking into each other. Every 100 steps the model
evaluates $J$ on the whole training set and stores it in \texttt{history\_},
which is how we check that the objective really decreases (Problem~3a); with
$2\cdot10^5$ steps that is 2{,}000 full-data evaluations, and because they
happen inside \texttt{fit} they are part of the fit time reported in
Table~\ref{tab:timing}. \texttt{decision\_function} returns the signed margin
$\mathbf{w}^{\top}\mathbf{x}+b$, which is also what the ROC curves and the AUC
use; a score of exactly zero is predicted as the positive class, as in
scikit-learn.

\subsection{Random forest from our own trees (bonus)}
\texttt{src/random\_forest.py} grows $B$ independent copies of our tree, each
on a bootstrap sample of the training rows and each considering only a random
subset of features at every node (\texttt{max\_features}, implemented inside
the tree so that constant columns in a node are skipped rather than wasting
the draw). Classification averages the trees' \texttt{predict\_proba} (soft
voting) and takes the argmax; regression averages the predictions.
\texttt{staged\_predict} returns the prediction of the first $1,2,\dots,B$
trees from a single running sum, which is what Figure~\ref{fig:forest} plots.

\section{Experimental Setup}
\label{sec:setup}
\textbf{Protocol.} Every model sees the same 70/15/15 training, validation
and test split (seed 42). The trees' settings are chosen on the validation
split only (Table~\ref{tab:tree-study}); the SVM uses $\lambda=10^{-4}$ and
200{,}000 Pegasos updates, and its scikit-learn twins the same $\lambda$. The
Task~B threshold is the training median price, and the test split is scored
once, with the final settings. Each model is also cross-validated
on the training split (5 folds, stratified by tier for Task~B; mean $\pm$
std), and each headline test metric has a 95\% percentile bootstrap interval
(1{,}000 resamples of the test rows).

\textbf{Metrics.} Task~A: RMSE, MAE and $R^2$ on $\log(\text{price})$, the
scale the models are trained on, and the median absolute percentage error of
the price in AZN, which is easier to read. Task~B: precision, recall and F1
for the premium class, accuracy, ROC-AUC from each model's continuous scores,
and the confusion matrix. Every one of these, together with the stratified
$k$-fold splitter and the bootstrap, is implemented in
\texttt{src/evaluate.py} in NumPy and checked against scikit-learn's values in
\texttt{tests/test\_evaluate.py}.

\textbf{Task~A target for the trees.} A tree can only follow the floor area
in steps, so both trees are fitted on $\log(\text{price}/\text{area})$ (area
in m$^2$, or in sot for land plots) and the listing's own $\log(\text{area})$
is added back to the prediction. The area is an ordinary feature of the
listing; \texttt{unit\_price} and \texttt{total\_price} are dropped from
every feature matrix and price per m$^2$ is never an input, only the training
target. Because the area is added back exactly, the residuals are the same as
on $\log(\text{price})$, so every Task~A metric, $R^2$ included, is computed
on $\log(\text{price})$. The tree settings are tuned on the validation split
(\texttt{src/tree\_experiments.py}, Table~\ref{tab:tree-study}) and the same
settings are used in every table.

\textbf{Baselines.} scikit-learn's \texttt{DecisionTreeRegressor} and
\texttt{DecisionTreeClassifier} get our tree's criterion, depth and leaf
settings. Against our Pegasos SVM we fit \texttt{SGDClassifier} (hinge loss,
$\alpha=\lambda$, whole epochs with at least as many updates as Pegasos) and
\texttt{LinearSVC} ($C=1/(\lambda n)$ for $n$ training rows), each in a
\texttt{Pipeline} after \texttt{StandardScaler}. All three target the same
$\lambda$-regularised hinge loss on features standardised with training
statistics only, but none is run to the exact optimum: Pegasos and SGD stop
after a fixed number of updates, \texttt{LinearSVC} after 20{,}000
iterations (with a convergence warning), and each handles the intercept its
own way. Trivial floors predict the training mean (Task~A) and the majority
class (Task~B).

\textbf{Analysis and timing.} Permutation importance (on the validation
split, for every model) and the trees' impurity decrease treat all one-hot
columns of a feature as one; for the linear SVMs we show the largest
individual weights on the standardised columns. Test errors are broken down
by fixed price bands, property category, the number of training listings in
the listing's district and the distance from the tier threshold, and we list
the listings our tree misprices most.
Fit and predict times are wall-clock seconds measured in one run on one
laptop (Intel Core i5-14450HX); they are comparable with each other, not
across machines. Fit time against training-set size is the median of three
fits on nested subsamples of $1/8$, $1/4$, $1/2$ and all of the training
split.

\textbf{Reproducibility.} All randomness is seeded, and
\texttt{python -m src.run\_all} regenerates every number, table and figure
in this report from the raw data; a test checks that two runs give identical
numbers apart from the timings. A CI check
(\texttt{tools/check\_golden\_rule.py}) parses the imports of
\texttt{decision\_tree.py}, \texttt{random\_forest.py} and \texttt{svm.py},
and of our own modules they import, and fails the build if a model or
optimiser from any library reaches them.

\section{Results}
\label{sec:results}
Every number below is produced by one seeded run of
\texttt{python -m src.run\_all}. Table~\ref{tab:results} is the headline
comparison on the held-out test split; Tables~\ref{tab:task-a}
and~\ref{tab:task-b} add the cross-validated scores and the remaining
metrics, Table~\ref{tab:timing} the fit and predict times, and
Table~\ref{tab:tree-study} the tree study behind the settings.

\subsection{Choosing the trees on the validation split}
Both trees are tuned over the same grid --- \texttt{max\_depth} $\in\{2, 4, 6,
8, 10, 12, 15, 20, \text{unlimited}\}$ crossed with
\texttt{min\_samples\_leaf} $\in\{1, 5, 10, 20, 50\}$ --- scored on the
validation split only.

For Task~B (Figure~\ref{fig:tree-study}, left) the best entropy tree
(depth~20, 1-row leaves, validation F1 $0.9179$) and the best Gini tree
(depth~15, 1-row leaves, $0.9176$) are indistinguishable: the gap between the
criteria is a tenth of the $\pm0.004$ spread we later see across
cross-validation folds. We keep entropy because the rule was to take the best
validation score, but the honest reading is that the criterion does not
matter here. Depth does: validation F1 climbs from $0.814$ at depth~2 to
$0.911$ at depth~10 and then flattens.

Task~A shows the bias--variance trade-off much more sharply, and shows why
depth and leaf size must be tuned \emph{together} (Figure~\ref{fig:tree-grid}).
With 1-row leaves, validation RMSE falls to $0.357$ at depth~8 and then
deteriorates to $0.398$ as the tree is grown out --- textbook overfitting,
visible as the widening gap between the dashed training curve and the solid
validation curve in Figure~\ref{fig:tree-study} (right). With 20-row leaves
the same trees keep improving as they deepen, down to $0.335$ at unlimited
depth, which is what the grid selects. Tuning depth first at 1-row leaves and
the leaf size afterwards would have stopped in a clearly worse corner of the
grid. The learning curves of the two final trees (Figure~\ref{fig:learning}) show
that they sit in different regimes. The Task~A tree, with 20-row leaves, ends
at a training RMSE of $0.288$ against a validation RMSE of $0.335$: a modest
gap, and the validation curve is still falling ($0.345$ at half the training
data), so it is limited by data rather than over-fitted. The Task~B tree, with
1-row leaves, memorises its training split (F1 $\approx 0.99$) while reaching
$0.918$ on validation --- a high-variance model that validation still prefers,
and exactly the kind of learner the random forest in
Section~\ref{subsec:forest} averages to good effect.

\subsection{Task A: predicting the price}
Our tree reaches a test RMSE of $0.332\ [0.317, 0.349]$ on
$\log(\text{price})$, $R^2 = 0.829$, MAE $0.185$, and a median absolute
percentage error of $10.3\,\%$ in AZN: half the test listings are priced
within about a tenth of their asking price. Five-fold cross-validation on the
training split gives $0.334 \pm 0.013$, so the test split is not a lucky draw.
The trivial floor that predicts the training mean scores RMSE $0.803$ and a
median error of $41.7\,\%$, and the bootstrap intervals of the two are far
apart. scikit-learn's \texttt{DecisionTreeRegressor} with our settings lands on
the same numbers to three decimals (Table~\ref{tab:task-a}).
Figure~\ref{fig:pred-actual} shows where the remaining error sits: the cloud
hugs the diagonal in the middle of the price range and bends towards the mean
at both ends.

\subsection{Task B: premium vs standard}
Our tree separates the two tiers with F1 $=0.921\ [0.916, 0.927]$, accuracy
$0.921$ and AUC $0.927$, a hair above scikit-learn's $0.919$ and $0.925$ with
identical settings. Our Pegasos SVM reaches F1 $=0.867\ [0.860, 0.874]$ and
AUC $0.930$: better than \texttt{SGDClassifier} on the same objective and
update budget (F1 $0.849$, AUC $0.919$), below \texttt{LinearSVC}
(F1 $0.893$, AUC $0.954$). Section~\ref{sec:discussion} argues that this last
gap is an optimisation budget rather than a different model.

Two details in Table~\ref{tab:task-b} are worth pointing out. First, our SVM's
AUC is marginally \emph{higher} than our tree's even though its F1 is five
points lower: the margin ranks listings well (Figure~\ref{fig:roc}), but the
sign of the margin is not where F1 peaks, while the tree's leaf-frequency
probabilities are coarse, which caps its AUC rather than its accuracy. Second,
the trivial floor has precision, recall and F1 of exactly zero, because the
majority class in the training split is ``standard'' and it never predicts
``premium''. Figure~\ref{fig:confusion} shows how the errors of each
classifier divide between false and missed premiums: our tree is almost
symmetric (precision $0.923$, recall $0.920$), while \texttt{SGDClassifier}
misses premiums more often than it invents them (precision $0.862$, recall
$0.836$).

\subsection{Our implementations against scikit-learn}
Given the same hyperparameters and the same rows, our trees and
scikit-learn's agree on $99.2\,\%$ (Task~A) and $98.4\,\%$ (Task~B) of the
test predictions, reach the same depth (29 and 20) and have almost the same
size: an identical 1{,}724 leaves on Task~A, and 2{,}690 against 2{,}685 on
Task~B. They are nevertheless not identical node by node, and the study
(\texttt{tree\_experiments.compare\_with\_sklearn}) separates the two reasons.
Equally good splits are broken differently --- we always take the lowest
feature index and the smallest threshold, scikit-learn draws features in a
random order, so a seed search is enough to reproduce our tree exactly on
tie-free problems --- and scikit-learn casts $X$ to \texttt{float32} before
fitting, which merges feature values that differ only in the seventh digit,
something our \texttt{float64} search keeps apart. On the test split the
largest single disagreement is $0.156$ in $\log(\text{price})$, i.e.\ one
listing landing in a neighbouring leaf, and the mean absolute difference
rounds to zero.

The price of writing the models ourselves is speed, not accuracy
(Table~\ref{tab:timing}). On the full training split our tree fits in
$3.3$~s against $0.24$~s for scikit-learn's (Task~A), and our SVM in
$3.2$~s against $0.086$~s for \texttt{SGDClassifier} --- although
\texttt{LinearSVC}, which actually solves the problem, takes $3.6$~s.
Prediction is 5--7~ms against 2~ms. Figure~\ref{fig:scaling} shows the shape
of the gap: eight times more training rows costs our tree about fourteen times
more time ($0.21$~s at 5{,}552 rows, $3.06$~s at 44{,}423), the same mildly
super-linear $O(nd\log n)$ behaviour as scikit-learn's curve, shifted up by a
constant factor of roughly thirteen. That factor is interpreter overhead per
node, not a worse algorithm: scikit-learn's splitter is compiled Cython. Our
SVM takes a fixed 200{,}000 steps whatever the data size, so its fit time
grows only because the 2{,}000 full-training-set evaluations of the objective
that it records for the convergence check scale with $n$.

\subsection{What the models rely on}
Permutation importance on the validation split (Figure~\ref{fig:importance})
puts the floor area first for every model on both tasks: shuffling it costs
our Task~A tree $0.594$ RMSE points, far more than the distance to the city
centre ($0.431$), the land area ($0.319$), the longitude ($0.119$) or the
property category ($0.113$). The trees' impurity-decrease importances
(Figure~\ref{fig:impurity}) rank the same features slightly differently ---
area $0.41$, category $0.21$, distance $0.20$ --- which is the usual
disagreement between ``how much the split gained during training'' and ``how
much the prediction degrades without it''. On Task~A our tree's and
scikit-learn's importances agree to three decimals.

The linear SVMs tell a compatible story through their weights
(Figure~\ref{fig:svm-weights}): on standardised columns our Pegasos SVM gives
the floor area a weight of $+21.0$, land area $+2.8$ and the land category
$+2.7$, against $-2.7$ for the distance to the centre and $-3.9$, $-3.1$ for
the cheapest outlying districts. \texttt{LinearSVC} orders the same features
with the same signs but with weights roughly $2.5$ times smaller --- a wider
margin for the same $\lambda$, the signature of a solver that has reached the
optimum of the objective we all share.

\subsection{Where the errors are}
Three breakdowns explain almost all of the residual error.

\emph{Price level} (Figure~\ref{fig:errors-a}, left). The median relative
error is $8.2\,\%$ in the 100--200k~AZN band and $8.7\,\%$ in 200--300k, but
$16.7\,\%$ below 100k and $38.7\,\%$ above 1~M~AZN. The bias column shows the
mechanism: $+19\,\%$ on the cheapest band and $-36\,\%$ on the most expensive
one. A tree predicts a leaf mean, so it cannot extrapolate past the range it
has seen and shrinks both tails towards the middle.

\emph{Property type and district coverage} (Table~\ref{tab:errors}). Flats ---
new build and old stock, $76\,\%$ of the test split --- are priced to within
$8.1\,\%$, houses to $21.1\,\%$, commercial units to $37.0\,\%$ and land plots
to $40.0\,\%$. Likewise, listings in districts with at least 300 training
listings have a median error of $9.5\,\%$ against $26.3\,\%$ in districts with
fewer than 30, where the one-hot column has been pooled into ``other''; the
Task~B error rate rises from $7.5\,\%$ to $10.9\,\%$ for our tree and from
$13.0\,\%$ to $24.8\,\%$ for our SVM over the same groups. The listings our
tree misprices most (Table~\ref{tab:worst}) are almost all land plots, where
the price per sot ranges over four orders of magnitude between a suburban
garden plot and development land.

\emph{Distance from the tier threshold} (Figure~\ref{fig:errors-b}). Within
$5\,\%$ of the 209{,}000~AZN threshold our tree misclassifies $37\,\%$ of
listings and our SVM $49\,\%$; beyond half a log unit the error rates fall to
$3.0\,\%$ and $5.5\,\%$. Almost every Task~B mistake is a listing that is
genuinely borderline, which is what a median split on a continuous quantity
guarantees.

\subsection{Bonus: a forest of our own trees}
\label{subsec:forest}
Averaging 100 of our own trees, each on a bootstrap sample and with a random
subset of features per node, cuts the Task~A test RMSE from $0.332$ to
$0.291$ and lifts the Task~B F1 from $0.921$ to $0.944$
(Figure~\ref{fig:forest}). Both beat scikit-learn's \texttt{RandomForest} with
the same settings ($0.296$ and $0.940$). The gain arrives quickly --- ten
trees already give RMSE $0.301$ and F1 $0.937$, twenty give $0.296$ and
$0.941$ --- and the curves are nearly flat beyond about forty trees
($0.293$ and $0.942$), which is the variance reduction an ensemble of
decorrelated high-variance learners is supposed to deliver. The cost is again
constant-factor speed: 76~s (Task~B) and 186~s (Task~A) against 3.0~s and
6.8~s for scikit-learn.

\generatedtable{tables/headline.tex}
\generatedtable{tables/task_a.tex}
\generatedtable{tables/task_b.tex}
\generatedtable{tables/timing.tex}
\generatedtable{tables/tree_study.tex}

\begin{figure*}[t]
  \centering
  \generatedfigure{figures/tree_b_depth_gini_entropy.png}{0.48\textwidth}\hfill
  \generatedfigure{figures/tree_a_depth.png}{0.48\textwidth}
  \caption{Tree study on the validation split. Left: Task~B F1 against
    max\_depth for Gini and entropy, each at its own tuned
    min\_samples\_leaf. Right: Task~A RMSE of $\log(\text{price})$ against
    max\_depth with 1-row leaves and, if different, with the tuned leaf
    size. Dashed: training split, solid: validation split.}
  \label{fig:tree-study}
\end{figure*}
\begin{figure}[t]
  \centering
  \generatedfigure{figures/tree_a_grid.png}{\linewidth}
  \caption{Task~A validation RMSE of $\log(\text{price})$ over max\_depth
    $\times$ min\_samples\_leaf; the red frame marks the chosen setting.}
  \label{fig:tree-grid}
\end{figure}
\begin{figure*}[t]
  \centering
  \generatedfigure{figures/tree_b_learning_curve.png}{0.48\textwidth}\hfill
  \generatedfigure{figures/tree_a_learning_curve.png}{0.48\textwidth}
  \caption{Learning curves of the final trees on growing random subsets of
    the training split (left: Task~B F1, right: Task~A RMSE of
    $\log(\text{price})$). Dashed: the score on the training rows used, solid:
    on the fixed validation split.}
  \label{fig:learning}
\end{figure*}
\begin{figure}[t]
  \centering
  \generatedfigure{figures/eval_pred_vs_actual.pdf}{\linewidth}
  \caption{Task~A: predicted vs actual $\log(\text{price})$ on the test
    split; the dashed line is a perfect prediction.}
  \label{fig:pred-actual}
\end{figure}
\begin{figure*}[t]
  \centering
  \generatedfigure{figures/eval_confusion_matrices.pdf}{\textwidth}
  \caption{Task~B confusion matrices on the test split (rows: actual
    tier, columns: predicted tier).}
  \label{fig:confusion}
\end{figure*}
\begin{figure}[t]
  \centering
  \generatedfigure{figures/eval_roc_curves.pdf}{\linewidth}
  \caption{Task~B ROC curves on the test split, from each model's
    continuous scores (the SVMs' margins, the trees' predicted probabilities).}
  \label{fig:roc}
\end{figure}
\begin{figure}[t]
  \centering
  \generatedfigure{figures/eval_fit_time_scaling.pdf}{\linewidth}
  \caption{Fit time against training-set size (median of three fits on
    nested subsamples of the training split; log--log axes). SGD's epochs
    and \texttt{LinearSVC}'s $C$ are re-matched at each size;
    Table~\ref{tab:timing} notes anything else our SVM's fit time
    includes.}
  \label{fig:scaling}
\end{figure}
\begin{figure*}[t]
  \centering
  \generatedfigure{figures/eval_feature_importance.pdf}{0.8\textwidth}
  \caption{Permutation importance on the validation split: how much the
    score drops (RMSE on $\log(\text{price})$ rises for Task~A, F1 falls
    for Task~B) when all one-hot columns of a feature are shuffled
    together; mean of three shuffles, largest twelve features per model.}
  \label{fig:importance}
\end{figure*}
\begin{figure}[t]
  \centering
  \generatedfigure{figures/eval_tree_impurity_importance.pdf}{\linewidth}
  \caption{Impurity-decrease importance of the trees on the training
    split, as a share of the total decrease, with the one-hot columns of a
    feature summed. For Task~A the log area is added back outside the
    tree, so the area's importance here understates how much it drives the
    price.}
  \label{fig:impurity}
\end{figure}
\begin{figure*}[t]
  \centering
  \generatedfigure{figures/eval_svm_weights.pdf}{0.85\textwidth}
  \caption{The twelve largest weights (by absolute value) of each linear
    SVM, on standardised columns, one-hot columns shown separately; blue
    pushes a listing towards premium, red towards standard. Each panel has
    its own scale.}
  \label{fig:svm-weights}
\end{figure*}
\begin{figure*}[t]
  \centering
  \begin{minipage}[b]{0.56\textwidth}
    \centering
    \generatedfigure{figures/eval_error_by_price_band.pdf}{\linewidth}
  \end{minipage}\hfill
  \begin{minipage}[b]{0.42\textwidth}
    \centering
    \generatedfigure{figures/eval_error_by_district_coverage.pdf}{\linewidth}
  \end{minipage}
  \caption{Task~A test error by price band (left) and by the number of
    training listings in the listing's district (right): median of
    $|$predicted $-$ actual$|$ / actual price; $n$ = test listings in the
    group.}
  \label{fig:errors-a}
\end{figure*}
\begin{figure}[t]
  \centering
  \generatedfigure{figures/eval_error_by_threshold_margin.pdf}{\linewidth}
  \caption{Task~B test error rate by distance from the tier threshold,
    $|\log(\text{price}/\text{threshold})|$; $n$ = test listings in the
    group.}
  \label{fig:errors-b}
\end{figure}
\begin{figure*}[t]
  \centering
  \generatedfigure{figures/forest_b_n_trees.png}{0.48\textwidth}\hfill
  \generatedfigure{figures/forest_a_n_trees.png}{0.48\textwidth}
  \caption{Random forest of our own trees (Problem~9): validation and test
    score as trees are added, for Task~B (F1, left) and Task~A (RMSE of
    $\log(\text{price})$, right), against our single tuned tree (grey,
    dashed) and scikit-learn's RandomForest with the same settings (red,
    dotted), both on the test split.}
  \label{fig:forest}
\end{figure*}


\generatedtable{tables/errors.tex}
\generatedtable{tables/worst.tex}

\section{Discussion and Limitations}
\label{sec:discussion}
\textbf{Trees, lines and the shape of the problem.} The tree wins both tasks
by a wide margin over the linear SVM (F1 $0.921$ against $0.867$) because
price in this market is built from interactions: the value of a square metre
depends on the district, the category and the distance to the centre at the
same time. An axis-aligned tree represents such a product structure with a few
nested splits; a linear model on one-hot columns can only add the effects up,
and the one interaction it is given --- the per-m$^2$ target --- is the one we
hand-engineered for the tree. The SVM nevertheless ranks listings almost as
well as the tree (AUC $0.930$ against $0.927$), which says that the premium
tier is close to linearly separable in these features and that most of the F1
gap is about \emph{where} the boundary sits rather than its shape.

\textbf{Why our SVM trails \texttt{LinearSVC}.} All three linear models
minimise the same $\lambda$-regularised hinge loss on the same standardised
features, so the $2.6$-point F1 gap is about optimisation, not formulation.
Pegasos takes 200{,}000 stochastic steps, about $4.5$ passes over 44{,}423
rows, and \texttt{SGDClassifier} is given at least the same budget;
\texttt{LinearSVC} runs coordinate descent on the dual for up to 20{,}000
iterations and still raises a \texttt{ConvergenceWarning} on all six of its
fits (the full training split and the five cross-validation folds). The weight
magnitudes confirm the diagnosis: our solution's weights are about $2.5$ times
larger than \texttt{LinearSVC}'s, i.e.\ our margin is narrower than the
optimum of the objective we set out to minimise. More steps, averaging the
iterates, or a mini-batch sub-gradient would close most of that gap.

\textbf{Bias and variance.} The depth sweeps make the trade-off concrete: a
depth-2 tree is pure bias (validation RMSE $0.470$, F1 $0.814$), a fully grown
tree with 1-row leaves is pure variance ($0.398$), and the useful models sit
between the two --- in Task~A's case only once the leaf size is raised to 20,
after which depth stops hurting at all. The random forest attacks the same
variance from the other direction and buys a further $12\,\%$ of RMSE.

\textbf{What we would not deploy.} The model is a good estimator of a typical
Baku flat and a poor estimator of everything else: luxury listings above
1~M~AZN are under-priced by about a third on average, land plots and
commercial units carry a $37$--$40\,\%$ median error, and thinly covered
districts nearly triple the error of well-covered ones. A deployed version would need either separate
models per property type or a target that is less dominated by flats.

\textbf{Limitations.} We list the ones that affect how the numbers should be
read.
\begin{enumerate}
  \item \emph{The SVM is not tuned.} $\lambda = 10^{-4}$ and 200{,}000 steps
    were fixed before training and no script sweeps them, so Problem~4b is
    only partly addressed: we show that the objective decreases, but not how
    the margin and the validation score trade off against $\lambda$. The
    \texttt{LinearSVC} comparison suggests a better solution exists at the
    very same $\lambda$.
  \item \emph{The SVM is linear.} We did not implement the non-linear
    extension (kernelised Pegasos or random Fourier features, Problem~3d), so
    the SVM cannot represent the interactions the tree exploits.
  \item \emph{Margin diagnostics are missing.} The pipeline reports the
    largest standardised weights but not $\lVert\mathbf{w}\rVert$, the margin
    width $2/\lVert\mathbf{w}\rVert$ or the identity of the support vectors,
    which Problem~3c asks for.
  \item \emph{One split, one seed.} Tuning used a single validation split and
    the test split was scored once. The bootstrap intervals cover test
    sampling noise only, not the variability of the split itself; repeated
    splits or nested cross-validation would widen them.
  \item \emph{Asking prices, not transactions.} The target is what sellers
    ask on a 2024--25 snapshot in which $99.6\,\%$ of listings are in Baku.
    Nothing here estimates what properties actually sold for, and the model
    should not be read as a valuation.
  \item \emph{Spatial generalisation is untested.} Districts are coarse (142
    names, 92 kept and the rest pooled into ``other'') and we never evaluate
    on a district that is absent from training; a spatially blocked split
    would be a harder and more honest test.
  \item \emph{Cleaning involves judgement.} The outlier rules of
    Section~\ref{sec:data} are fixed domain thresholds, not learned ones, but
    they are still our choices, and they remove $0.13\,\%$ of rows; the re-post key is deliberately strict, so some
    duplicate properties certainly survive it.
  \item \emph{Unused signal.} The free-text description and the attribute
    string are used only to recover a missing room count, and the building
    type is unusable ($99.8\,\%$ missing).
  \item \emph{Speed.} Our implementations are 13--40 times slower to fit than
    scikit-learn's compiled equivalents. That is comfortable at 44{,}423 rows
    and would not be at ten million.
\end{enumerate}

\section{Conclusion}
Two models written from scratch in NumPy --- a CART tree with a vectorised
prefix-sum split search and a Pegasos sub-gradient SVM --- reproduce
scikit-learn's behaviour on a real, messy Azerbaijani property dataset: the
tree matches \texttt{DecisionTreeRegressor} to three decimals on Task~A
(RMSE $0.332$ on $\log(\text{price})$, median error $10.3\,\%$) and agrees
with it on $99\,\%$ of individual predictions, while on Task~B it reaches F1
$0.921$ against scikit-learn's $0.919$, and our Pegasos SVM reaches F1
$0.867$, between \texttt{SGDClassifier} and \texttt{LinearSVC}. Writing them
ourselves cost a constant factor in speed and bought a precise understanding
of where the remaining error lives: in the price tails, in land and commercial
listings, in thinly covered districts, and within a few per cent of the tier
threshold. The clearest next steps follow directly from the limitations above
--- sweep $\lambda$ and the step budget for the SVM and add a random-Fourier
feature map, give the tails their own treatment (a quantile loss, or separate
models per property type), and add neighbourhood price features that let a
sparse district borrow strength from its neighbours.

\section*{Contributions}
\textbf{Usubov, Tural} --- the from-scratch decision tree
(\texttt{src/decision\_tree.py}), its tuning, bias--variance study and
node-by-node comparison with scikit-learn
(\texttt{src/tree\_experiments.py}), the random-forest bonus
(\texttt{src/random\_forest.py}) and their tests (PR~\#6).
\textbf{Huseynli, Ali} --- the data pipeline: cleaning, de-duplication,
outlier rules, feature engineering, the leakage-free split and the EDA
figures (\texttt{src/data\_prep.py}, \texttt{src/eda.py},
\texttt{eda.ipynb}; PRs~\#10, \#11), and contributions to the tree work.
\textbf{Ibrahimli, Polad} --- the from-scratch Pegasos SVM
(\texttt{src/svm.py}; PR~\#7).
\textbf{Eyyubov, Samur} --- evaluation and reproduction: the metric,
cross-validation and bootstrap module, the matched scikit-learn baselines,
the error analysis and plots, the one-command runner and the repository
workflow (\texttt{src/evaluate.py}, \texttt{src/baselines.py},
\texttt{src/error\_analysis.py}, \texttt{src/plots.py},
\texttt{src/run\_all.py}, CI; PRs~\#1, \#3, \#5).
\texttt{contribution\_report.md} in the repository root gives the same split
with the effort percentages.

\section*{Tools and Acknowledgements}
The dataset is the \emph{Bina.az Sale Project} dump published on Kaggle by
Şəhriyar Məmmədli~\cite{dataset}; we use it
under its Kaggle licence and did not re-scrape or redistribute it. Our own
models use only NumPy~\cite{numpy}; pandas, Matplotlib and scikit-learn are
used for data handling, figures and the comparison baselines~\cite{sklearn}.
We used the Claude AI assistant (Anthropic) while working on this project: for
the repository and CI scaffolding, for review of our modules, and for drafting
and copy-editing this report. Every design decision, every number and every
claim here was checked by us against the code and the pipeline output, and all
model code in \texttt{src/decision\_tree.py}, \texttt{src/svm.py} and
\texttt{src/random\_forest.py} is our own NumPy implementation, which
\texttt{tools/check\_golden\_rule.py} enforces in CI.

% =====================================================================
\begin{thebibliography}{9}
\bibitem{cart}
L.~Breiman, J.~H. Friedman, R.~A. Olshen, and C.~J. Stone,
\emph{Classification and Regression Trees}. Belmont, CA: Wadsworth, 1984.

\bibitem{cortes}
C.~Cortes and V.~Vapnik, ``Support-vector networks,''
\emph{Machine Learning}, vol.~20, no.~3, pp.~273--297, 1995.

\bibitem{pegasos}
S.~Shalev-Shwartz, Y.~Singer, N.~Srebro, and A.~Cotter,
``Pegasos: primal estimated sub-gradient solver for SVM,''
\emph{Mathematical Programming}, vol.~127, no.~1, pp.~3--30, 2011.

\bibitem{sklearn}
F.~Pedregosa \emph{et al.}, ``Scikit-learn: machine learning in Python,''
\emph{Journal of Machine Learning Research}, vol.~12, pp.~2825--2830, 2011.

\bibitem{numpy}
C.~R. Harris \emph{et al.}, ``Array programming with NumPy,''
\emph{Nature}, vol.~585, pp.~357--362, 2020.

\bibitem{dataset}
Ş.~Məmmədli, ``Bina.az Sale Project,'' Kaggle dataset,
2025. [Online]. Available:
\url{https://www.kaggle.com/datasets/sehriyarmemmedli/binaaz-sale-project}
\end{thebibliography}

\end{document}
